\begin{helper}
Let $G$ be a finite $\Delta$-regular simple graph.  Let $r$ be a vertex and
let $X\subseteq N_G(r)$ be such that no two distinct vertices of $X$ have a
common neighbor outside $N_G[r]$.  Let $Q\subseteq X$ be an independent set,
and let $P\subseteq Q$ have size two.  Let $\ell$ be a real-valued function
on finite vertex sets.  Suppose that, for every finite set $P'\subseteq X$
of size two that is independent in $G$, there are two distinct vertices
$u,v\in P'$ such that
\[
\ell(P')=\frac{|N_G(u)\cap N_G(v)|}{\Delta}.
\]
Then
\[
0\le \ell(P),\qquad 0<1-\ell(P),\qquad 0<2-\ell(P).
\]
\end{helper}
\begin{proof}
Since $P\subseteq Q\subseteq X$, the set $P$ is a subset of $X$.  Also $P$
is independent: if $a,b\in P$ are distinct, then $a,b\in Q$, and the
independence of $Q$ gives that $a$ and $b$ are not adjacent.  Therefore the
representation hypothesis for the function $\ell$ applies to this particular
set $P$.  Choose the two distinct vertices $u,v\in P$ supplied by that
hypothesis, so that
\[
\ell(P)=\frac{|N_G(u)\cap N_G(v)|}{\Delta}.
\]
Because $P\subseteq Q\subseteq X$, both $u$ and $v$ belong to $X$, and hence
both are adjacent to $r$.  Thus $\{u,v\}\subseteq N_G(r)$.  Because $G$ is
$\Delta$-regular, $|N_G(r)|=\Delta$, so in particular $\Delta\ge 2$.

We claim that every common neighbor of $u$ and $v$ lies in
\[
\{r\}\cup\bigl(N_G(r)\setminus\{u,v\}\bigr).
\]
Indeed, let $w\in N_G(u)\cap N_G(v)$.  If $w=r$ there is nothing to prove.
If $w\ne r$, the clean-neighborhood hypothesis applied to the distinct
vertices $u,v\in X$ rules out the possibility that $w$ is not adjacent to
$r$.  Hence $w\in N_G(r)$.  Also $w\ne u$ and $w\ne v$, since a simple graph
has no loops and $w$ is adjacent to $u$ and to $v$.  This proves the claimed
containment.

The set on the right has cardinality $\Delta-1$: the set
$N_G(r)\setminus\{u,v\}$ has size $\Delta-2$, and $r\notin N_G(r)$, so adding
$r$ contributes one new vertex.  Therefore
\[
|N_G(u)\cap N_G(v)|\le \Delta-1<\Delta .
\]
Dividing by the positive real number $\Delta$ gives
$0\le \ell(P)<1$.  The two strict inequalities
$0<1-\ell(P)$ and $0<2-\ell(P)$ follow immediately.
\end{proof}
